% ==================================================================
% id:       seq-025
% topic:    סדרות
% subtopic: חסימות של סדרות
% solution: full
% ==================================================================
\textbf{תהיינה
$(a_n)_{n=1}^{\infty}$
ו-$(b_n)_{n=1}^{\infty}$
סדרות חסומות. הוכיחו או הפריכו:}
\begin{enumerate}
\item
הסדרה
$(a_n+b_n)$
חסומה.
\hint{השתמשו באי-שוויון המשולש.}
\begin{solbox}
הטענה נכונה. מכיוון שהסדרות חסומות, קיימים
$M_1,M_2>0$
כך ש-$|a_n|\le M_1$ ו-$|b_n|\le M_2$ לכל $n$.
לפי אי-שוויון המשולש, לכל $n$
$$|a_n+b_n|\le|a_n|+|b_n|\le M_1+M_2 ,$$
ולכן
$(a_n+b_n)$
חסומה.
\end{solbox}
\item
הסדרה
$(a_n\cdot b_n)$
חסומה.
\hint{היעזרו בכך ש-$|xy|=|x|\cdot|y|$.}
\begin{solbox}
הטענה נכונה. עם $M_1,M_2$ כמו בסעיף א, לכל $n$
$$|a_nb_n|=|a_n|\cdot|b_n|\le M_1M_2 ,$$
ולכן
$(a_nb_n)$
חסומה.
\end{solbox}
\item
אם
$b_n\neq0$
לכל $n$, אז הסדרה
$\lr{\frac{a_n}{b_n}}$
חסומה.
\hint{מה קורה למנה כשהמכנה מתקרב לאפס?}
\begin{solbox}
הטענה לא נכונה. דוגמה נגדית:
$a_n=1$ ו-$b_n=\frac1n$.
שתי הסדרות חסומות (כי $|a_n|=1$ ו-$0<b_n\le1$), ו-$b_n\neq0$ לכל $n$.
אבל
$\frac{a_n}{b_n}=n$,
והסדרה הזו לא חסומה מלעיל: לכל $M$ יש $n>M$ (תכונת ארכימדס).
\end{solbox}
\item
אם קיים
$c>0$
כך ש-$|b_n|\ge c$
לכל $n$, אז הסדרה
$\lr{\frac{a_n}{b_n}}$
חסומה.
\hint{מה ההבדל בין התנאי כאן לתנאי בסעיף הקודם? מה הוא מונע מהמכנה לעשות?}
\begin{solbox}
הטענה נכונה. עם $M_1$ כמו בסעיף א, לכל $n$
$$\left|\frac{a_n}{b_n}\right|=\frac{|a_n|}{|b_n|}\le\frac{M_1}{c} ,$$
ולכן הסדרה חסומה. ההבדל מסעיף ג: כאן המכנה לא יכול להתקרב לאפס.
\end{solbox}
\end{enumerate}
